\begin{ex}
Tính giá trị của biểu thức $a + 25$ với $a = 15$.
\shortans{40}
\loigiai{
$15 + 25 = 40$.
}
\end{ex}

\begin{ex}
Tính giá trị của biểu thức $b - 18$ với $b = 50$.
\shortans{32}
\loigiai{
$50 - 18 = 32$.
}
\end{ex}

\begin{ex}
Tính giá trị của biểu thức $120 - c$ với $c = 45$.
\shortans{75}
\loigiai{
$120 - 45 = 75$.
}
\end{ex}

\begin{ex}
Tính giá trị của biểu thức $m \times 6$ với $m = 8$.
\shortans{48}
\loigiai{
$8 \times 6 = 48$.
}
\end{ex}

\begin{ex}
Tính giá trị của biểu thức $72 : n$ với $n = 9$.
\shortans{8}
\loigiai{
$72 : 9 = 8$.
}
\end{ex}

\begin{ex}
Tính giá trị của biểu thức $a + b$ với $a = 34$ và $b = 26$.
\shortans{60}
\loigiai{
$34 + 26 = 60$.
}
\end{ex}

\begin{ex}
Tính giá trị của biểu thức $a - b$ với $a = 85$ và $b = 38$.
\shortans{47}
\loigiai{
$85 - 38 = 47$.
}
\end{ex}

\begin{ex}
Tính giá trị của biểu thức $a \times b$ với $a = 12$ và $b = 5$.
\shortans{60}
\loigiai{
$12 \times 5 = 60$.
}
\end{ex}

\begin{ex}
Tính giá trị của biểu thức $a : b$ với $a = 96$ và $b = 4$.
\shortans{24}
\loigiai{
$96 : 4 = 24$.
}
\end{ex}

\begin{ex}
Tính giá trị của biểu thức $a + b + c$ với $a = 15, b = 25, c = 30$.
\shortans{70}
\loigiai{
$15 + 25 + 30 = 70$.
}
\end{ex}

\begin{ex}
Tính giá trị của biểu thức $a \times b \times c$ với $a = 2, b = 3, c = 5$.
\shortans{30}
\loigiai{
$2 \times 3 \times 5 = 30$.
}
\end{ex}

\begin{ex}
Tính giá trị của biểu thức $(a + b) \times c$ với $a = 14, b = 6, c = 4$.
\shortans{80}
\loigiai{
$(14 + 6) \times 4 = 20 \times 4 = 80$.
}
\end{ex}

\begin{ex}
Tính giá trị của biểu thức $(a - b) : c$ với $a = 50, b = 14, c = 6$.
\shortans{6}
\loigiai{
$(50 - 14) : 6 = 36 : 6 = 6$.
}
\end{ex}

\begin{ex}
Tính giá trị của biểu thức $a + b \times c$ với $a = 20, b = 5, c = 6$.
\shortans{50}
\loigiai{
Nhân trước: $20 + 30 = 50$.
}
\end{ex}

\begin{ex}
Tính giá trị của biểu thức $a - b : c$ với $a = 40, b = 24, c = 4$.
\shortans{34}
\loigiai{
Chia trước: $40 - 6 = 34$.
}
\end{ex}

\begin{ex}
Cho chu vi hình vuông có cạnh $a$ là $P = a \times 4$. Tính $P$ khi $a = 15$ cm.
\shortans{60}
\loigiai{
$P = 15 \times 4 = 60$ cm.
}
\end{ex}

\begin{ex}
Cho chu vi hình chữ nhật là $P = (a + b) \times 2$. Tính $P$ khi $a = 18$ cm, $b = 12$ cm.
\shortans{60}
\loigiai{
$P = (18 + 12) \times 2 = 60$ cm.
}
\end{ex}

\begin{ex}
Cho diện tích hình chữ nhật là $S = a \times b$. Tính $S$ khi $a = 15$ cm, $b = 6$ cm theo $\text{cm}^2$.
\shortans{90}
\loigiai{
$S = 15 \times 6 = 90\text{ cm}^2$.
}
\end{ex}

\begin{ex}
Cho diện tích hình vuông là $S = a \times a$. Tính $S$ khi $a = 8$ cm theo $\text{cm}^2$.
\shortans{64}
\loigiai{
$S = 8 \times 8 = 64\text{ cm}^2$.
}
\end{ex}

\begin{ex}
Tính giá trị của biểu thức $250 + m$ với $m = 750$.
\shortans{1000}
\loigiai{
$250 + 750 = 1000$.
}
\end{ex}

\begin{ex}
Tính giá trị của biểu thức $1000 - n$ với $n = 345$.
\shortans{655}
\loigiai{
$1000 - 345 = 655$.
}
\end{ex}

\begin{ex}
Tính giá trị của biểu thức $x \times 8 + y$ với $x = 10, y = 20$.
\shortans{100}
\loigiai{
$10 \times 8 + 20 = 100$.
}
\end{ex}

\begin{ex}
Với $x = 5$, giá trị của biểu thức $x \times x + 25$ bằng bao nhiêu?
\shortans{50}
\loigiai{
$5 \times 5 + 25 = 50$.
}
\end{ex}

\begin{ex}
Tính giá trị biểu thức $100 : (x - 5)$ với $x = 15$.
\shortans{10}
\loigiai{
$100 : (15 - 5) = 100 : 10 = 10$.
}
\end{ex}

\begin{ex}
Tính giá trị của biểu thức $(a + b) : 2$ với $a = 24, b = 36$.
\shortans{30}
\loigiai{
$(24 + 36) : 2 = 60 : 2 = 30$.
}
\end{ex}
